\scanpage{171}
\mspage{45}

(qu'on notera $\mathrm{Im}\,f$). Le morphisme
$Z\longrightarrow\mathrm{Im}\,f$ est un épimorphisme.

Cette dernière assertion \deleted{\unsure{revient}} provient de
celle-ci :

\uline{Lemme}. \unsure{5.?|5.3|5.7}. Soit \deleted{\unsure{$C$}} $\unsure{\mathcal M|M}$ une catégorie
\inserted{où} les noyaux de doubles flèches existent.
\inserted{Soit} $\unsure{g|f}:Z\longrightarrow Y$ une flèche de $\mathcal M$ \deleted{qui} \inserted{Alors $g$ est épi \unsure{ssi} $\unsure{\exists|]}$, $g$}
ne se factorise pas par un sous-objet \inserted{\uline{strict}} $Y_0$ de $Y$ différent de $Y$.
\deleted{Alors $g$ est un épimorphisme.}

\deleted{(Il suffit $=$ que $g$ ne se factorise pas
par un sous-objet \uline{strict} -- \unsure{notamment|autrement},
par les $\mathrm{Ker}(f_1,f_2)$ \ldots)}

On va faire une hypothèse supplémentaire
(pratiquement anodine) sur les
$\mathrm{Filt}^{\unsure{\pi|\varpi}}$\deleted{\unsure{,}} :

d) $\mathrm{Filt}^{\pi}$ est stable par limites inductives
indexées par des catégories d'indices $I$
telles que $\mathrm{card}^{\mathrm{fl}}(I)\leq\pi$\textsuperscript{*}. (Il suffit de
supposer la stabilité par objet initial, sommes
amalgamées simples, et $\underrightarrow{\mathrm{lim}}$ filtrantes indexées
par un ensemble ordonné $I$ de cardinal $\leq\pi$.)
% slanted
\margin{Cette condition sera nommée (\unsure{$\mathrm{d}_0$|d°}) par la suite,}

De plus, si $f:Z\longrightarrow X$ est une flèche
de $\mathcal M$ telle que \inserted{$Z$} \deleted{$\unsure{X|Z}$} $\in\mathrm{Filt}^{\pi}$, alors $\mathrm{Im}\,f\in\mathrm{Filt}^{\pi}$.

Cette dernière condition \inserted{sur $\mathrm{Im}\,f$} \deleted{résulte des}
\deleted{autres, quand les épimorphismes et les monomorphismes sont stricts,} se décompose en
deux

% archive sheet number: 170
